\documentclass{article}
\usepackage{url}
\usepackage{cite}
\usepackage{amsmath}
\usepackage{amssymb}

\title{INFO-F420: Project Proposal \\ Computing Smallest Convex Intersecting Polygon}
\author{Bataille Alex\\ Bouhdid Wassim}
\begin{document}
\maketitle
\section{Introduction}
The topic we chose is the computation of a smallest convex intersecting polygon for a set of objects in the plane \cite{Antoniadis_deBerg_Kisfaludi-Bak_Skarlatos_2025}. 
A polygon $C$ is an intersecting polygon for a set $\mathcal{O}$ of objects in $\mathbb{R}^2$ if $C$ intersects each object in $\mathcal{O}$, where the polygon includes 
its interior. The term ``smallest'' can refer to either the minimum-perimeter or the minimum-area convex intersecting polygon. 

The paper we chose \cite{Antoniadis_deBerg_Kisfaludi-Bak_Skarlatos_2025} presents an FPTAS for both problems for the case where $\mathcal{O}$ is a set of possibly intersecting convex polygons in the plane of total complexity $n$.
The exact segment algorithm is used to compute a smallest convex intersecting polygon for a set of line segments in the plane. The algorithm is based on a dynamic programming approach, and it is shown that the problem can be solved in polynomial time for a fixed number of segments.

\section{Project Objectives}
We want to understand the paper and the corresponding literature deeply enough, and grasp the main ideas of the algorithms, especially the exact segment 
algorithm and its analysis. 
This would be done by making an interactive webpage that allows the user to play with the different constructions and objects of the paper. 

\section{Work Distribution}
This is a group project, consisting of two members, Alex and Wassim. The way we will distribute the work is as follows: 
\begin{itemize}
    \item Wassim will be in charge of the understanding of the algorithm, focusing on the main ideas and the analysis.
    \item Alex will be in charge of the study of the literature, and the mathematical background of the problem.
\end{itemize}
Of course, both of us will work together and share our findings, such that we can both understand the project in its entirety. 

\bibliography{biblio}{}
\bibliographystyle{plain}

\end{document}